\documentclass[11pt,fontset=fandol]{ctexart}
\usepackage[a4paper,margin=23mm]{geometry}
\usepackage{amsmath,amssymb,mathtools,bm,mathrsfs,esint}
\usepackage{graphicx,longtable,booktabs,array,calc}
\usepackage[normalem]{ulem}
\usepackage{xcolor,hyperref}
\hypersetup{hidelinks}
\setcounter{secnumdepth}{-1}
\newcounter{none}
\providecommand{\tightlist}{\setlength{\itemsep}{0pt}\setlength{\parskip}{0pt}}
\providecommand{\passthrough}[1]{#1}
\usepackage{amssymb,mathrsfs,esint,bm,newunicodechar}
\allowdisplaybreaks
\setlength{\emergencystretch}{3em}
\providecommand{\llangle}{\langle\!\langle}
\providecommand{\rrangle}{\rangle\!\rangle}
\newunicodechar{∞}{\ensuremath{\infty}}
\newunicodechar{∈}{\ensuremath{\in}}
\newunicodechar{∑}{\ensuremath{\sum}}
\newunicodechar{∫}{\ensuremath{\int}}
\newunicodechar{≤}{\ensuremath{\leq}}
\newunicodechar{≥}{\ensuremath{\geq}}
\newunicodechar{≠}{\ensuremath{\neq}}
\newunicodechar{→}{\ensuremath{\to}}
\newunicodechar{⇒}{\ensuremath{\Rightarrow}}
\newunicodechar{α}{\ensuremath{\alpha}}
\newunicodechar{β}{\ensuremath{\beta}}
\newunicodechar{γ}{\ensuremath{\gamma}}
\newunicodechar{λ}{\ensuremath{\lambda}}
\newunicodechar{φ}{\ensuremath{\varphi}}
\newunicodechar{ε}{\ensuremath{\varepsilon}}
\newunicodechar{σ}{\ensuremath{\sigma}}
\newunicodechar{θ}{\ensuremath{\theta}}
\newunicodechar{δ}{\ensuremath{\delta}}
\newunicodechar{η}{\ensuremath{\eta}}
\newunicodechar{π}{\ensuremath{\pi}}
\newunicodechar{ξ}{\ensuremath{\xi}}
\newunicodechar{Σ}{\ensuremath{\Sigma}}
\newunicodechar{Δ}{\ensuremath{\Delta}}
\newunicodechar{∂}{\ensuremath{\partial}}
\newunicodechar{∀}{\ensuremath{\forall}}
\newunicodechar{∃}{\ensuremath{\exists}}
\newunicodechar{√}{\ensuremath{\surd}}
\newunicodechar{−}{\ensuremath{-}}
\newunicodechar{μ}{\ensuremath{\mu}}
\newunicodechar{ρ}{\ensuremath{\rho}}
\newunicodechar{τ}{\ensuremath{\tau}}
\newunicodechar{ω}{\ensuremath{\omega}}
\newunicodechar{Ω}{\ensuremath{\Omega}}
\newunicodechar{χ}{\ensuremath{\chi}}
\newunicodechar{ψ}{\ensuremath{\psi}}
\newunicodechar{Γ}{\ensuremath{\Gamma}}
\newunicodechar{Φ}{\ensuremath{\Phi}}
\newunicodechar{Λ}{\ensuremath{\Lambda}}
\newunicodechar{ℝ}{\ensuremath{\mathbb{R}}}
\newunicodechar{ℤ}{\ensuremath{\mathbb{Z}}}
\newunicodechar{ℕ}{\ensuremath{\mathbb{N}}}
\newunicodechar{△}{\ensuremath{\triangle}}
\newunicodechar{⊗}{\ensuremath{\otimes}}
\begin{document}
\section{2025年杭州电子科技大学601数学分析补充解析}\label{ux5e74ux676dux5ddeux7535ux5b50ux79d1ux6280ux5927ux5b66601ux6570ux5b66ux5206ux6790ux8865ux5145ux89e3ux6790}

本文件由助手按现有题面独立编写，全卷19个题号均保留原章节和编号。原参考解析为手写稿，现有OCR不足以可靠读取，因此本文件提供可复核的补充解答。题面缺失、条件不足或命题按原文不成立时，在对应题号明确说明，建议补充条件与原题条件分别列出。

\subsection{一、填空题}\label{ux4e00ux586bux7a7aux9898}

\subsubsection{1.
乘积比的极限}\label{ux4e58ux79efux6bd4ux7684ux6781ux9650}

\textbf{答案：}\(0\)。

令 \[
a_n=\frac{11\cdot12\cdots(n+10)}{1\cdot4\cdots(3n-2)}>0.
\] 则 \(a_{n+1}/a_n=(n+11)/(3n+1)\to1/3\)。存在 \(N\) 使 \(n\ge N\)
时此比值不超过 \(1/2\)，故 \(a_n\le a_N2^{-(n-N)}\to0\)。

\subsubsection{2. 积分方程}\label{ux79efux5206ux65b9ux7a0b}

\textbf{答案：}\(\displaystyle f(x)=x/2\)，\(x\ge0\)。

由原式 \(f(0)=0\)。记 \(F(x)=\int_0^xf(t)\,dt\)，积分存在的通常题意保证
\(F\) 连续，原式 \(f=\sqrt F\) 又保证 \(f\) 连续，故 \(F'=f\)。在
\(x>0\) 上，\(f(x)>0\)，可以对 \(f(x)^2=F(x)\) 求导，得到 \(2ff'=f\)，即
\(f'=1/2\)。所以 \(f(x)=x/2+C\)；由在零点连续知 \(C=0\)。代回有
\(\sqrt{\int_0^x(t/2)dt}=x/2\)，确实满足所有条件。

\subsubsection{3. 定积分}\label{ux5b9aux79efux5206}

\textbf{答案：}\(\displaystyle\frac{\alpha^{12}-(\alpha-1)^{12}}{12}\)。

因 \((x-1)(1-x)^{10}=(x-1)^{11}\)，直接积分： \[
\int_\alpha^{\alpha+1}(x-1)(1-x)^{10}\,dx
=\left[\frac{(x-1)^{12}}{12}\right]_\alpha^{\alpha+1}.
\]

\subsubsection{4.
含参数的极限与反常积分}\label{ux542bux53c2ux6570ux7684ux6781ux9650ux4e0eux53cdux5e38ux79efux5206}

按远程OCR所给的积分上限 \(c\)、被积函数 \(te^t\)
解答；本地OCR对此处字母识别有损。

\textbf{答案：}\(c=2\)。

左侧极限为 \(e^c\)，右侧由原函数 \((t-1)e^t\) 得 \[
\int_{-\infty}^cte^t\,dt=(c-1)e^c.
\] 因 \(e^c>0\)，原等式等价于 \(1=c-1\)。

\subsubsection{5.
四面体上的三重积分}\label{ux56dbux9762ux4f53ux4e0aux7684ux4e09ux91cdux79efux5206}

\textbf{答案：}\(1/8\)。

按原卷``四面体''的通常题意，区域为平面 \(x+y+z=1\) 与三个坐标平面围成的
\[
V=\{x,y,z\ge0:x+y+z\le1\}.
\]
原OCR写``三个坐标轴围成''，这里明确采用坐标平面这一必要几何含义。由对称性，
\[
\begin{aligned}
\iiint_V(x+y+z)\,dV
&=3\int_0^1\int_0^{1-x}\int_0^{1-x-y}x\,dz\,dy\,dx\\
&=\frac32\int_0^1x(1-x)^2\,dx=\frac18.
\end{aligned}
\]

\subsubsection{6.
椭圆上的曲线积分}\label{ux692dux5706ux4e0aux7684ux66f2ux7ebfux79efux5206}

本地OCR恢复的积分为 \(\oint_Lx\,dy\)，其中 \(L:x^2/4+y^2=1\)，逆时针。

\textbf{答案：}\(2\pi\)。

由 Green 公式，\(\oint_Lx\,dy=\iint_D1\,dx\,dy\)，等于半轴为 \(2,1\)
的椭圆面积 \(2\pi\)。也可取 \(x=2\cos t,y=\sin t\)，积分为
\(\int_0^{2\pi}2\cos^2t\,dt=2\pi\)。

\subsubsection{7.
隐函数拐点求参数，题干缺失}\label{ux9690ux51fdux6570ux62d0ux70b9ux6c42ux53c2ux6570ux9898ux5e72ux7f3aux5931}

原卷仅保留``隐函数中拐点，求参数''，没有隐函数方程、参数位置、所求点及正则条件，无法确定答案。

一般方法：对 \(F(x,y;\lambda)=0\)，在 \(F_y\ne0\) 的点有 \[
y'=-\frac{F_x}{F_y},\qquad
y''=-\frac{F_{xx}+2F_{xy}y'+F_{yy}(y')^2}{F_y}.
\] 把候选点代入 \(F=0\) 和 \(y''=0\)
可求候选参数，随后必须检验曲线两侧凹凸性是否改变。若
\(F_y=0\)，应另选局部参数或直接分析，不能套用上述分母公式。

\subsubsection{8.
圆锥体外表面的第一型曲面积分}\label{ux5706ux9525ux4f53ux5916ux8868ux9762ux7684ux7b2cux4e00ux578bux66f2ux9762ux79efux5206}

按``\(z=\sqrt{x^2+y^2}\) 与 \(z=1\) 所围的外表面''，\(S\)
包含圆锥侧面和顶面圆盘。

\textbf{答案：}\(\displaystyle\frac\pi2(1+\sqrt2)\)。

圆锥侧面上 \(z=r\)、\(0\le r\le1\)，\(dS=\sqrt2\,r\,dr\,d\theta\)，贡献
\[
\int_0^{2\pi}\int_0^1r^2\sqrt2r\,dr\,d\theta=\frac{\pi\sqrt2}{2}.
\] 顶面圆盘 \(z=1\) 的贡献为 \[
\int_0^{2\pi}\int_0^1r^2r\,dr\,d\theta=\frac\pi2.
\] 若原题的 \(S\) 专指圆锥侧面，则只取第一项
\(\pi\sqrt2/2\)；现有文字``所围（外表面）''支持包含顶面。

\subsection{二、计算题}\label{ux4e8cux8ba1ux7b97ux9898}

\subsubsection{1.
下半球面取上侧的第二型曲面积分}\label{ux4e0bux534aux7403ux9762ux53d6ux4e0aux4fa7ux7684ux7b2cux4e8cux578bux66f2ux9762ux79efux5206}

\textbf{答案：}\(\displaystyle-\frac{2\pi}{105}\)。

下半球为 \(z=-\sqrt{1-x^2-y^2}\)。因为取上侧，\(dx\,dy\)
的投影符号为正，因此 \[
\iint_Sx^2y^2z\,dx\,dy
=-\iint_{x^2+y^2\le1}x^2y^2\sqrt{1-x^2-y^2}\,dx\,dy.
\] 转为极坐标，结果为 \[
-\left(\int_0^{2\pi}\cos^2\theta\sin^2\theta\,d\theta\right)
\left(\int_0^1r^5\sqrt{1-r^2}\,dr\right)
=-\frac\pi4\cdot\frac8{105}=-\frac{2\pi}{105}.
\]

\subsubsection{2.
由曲线积分恒为零确定参数}\label{ux7531ux66f2ux7ebfux79efux5206ux6052ux4e3aux96f6ux786eux5b9aux53c2ux6570}

原题现有文字为 \(I=\int_LP\,dx+Q\,dy\equiv0\)， \[
P=\frac xy r^\alpha,\qquad Q=-\frac{x^2}{y^2}r^\alpha,\qquad r=\sqrt{x^2+y^2}.
\] \textbf{条件缺口：}没有给出 \(L\)
的定义，也没有明确``恒为零''是否指每一条闭曲线。若只要求某一条未知曲线的积分为零，不能唯一确定
\(\alpha\)。

\textbf{按标准补充条件解答：}假设对半平面 \(y>0\) 或 \(y<0\)
内的任意分段光滑闭曲线 \(L\) 都有 \(I=0\)，则 \(\boxed{\alpha=-1}\)。

必要性来自 \(Q_x=P_y\)。直接计算 \[
Q_x-P_y=-(\alpha+1)\frac{x}{y^2}r^\alpha.
\] 它在半平面内恒为零当且仅当 \(\alpha=-1\)。充分性也可直接验证：当
\(\alpha=-1\)， \[
d\left(\frac ry\right)
=\frac{x}{yr}\,dx-\frac{x^2}{y^2r}\,dy=P\,dx+Q\,dy,
\] 闭曲线上的积分必为零。曲线不得穿过
\(y=0\)，因为原积分系数在那里没有定义。

\subsubsection{3.
幂级数的收敛半径与和函数}\label{ux5e42ux7ea7ux6570ux7684ux6536ux655bux534aux5f84ux4e0eux548cux51fdux6570}

\textbf{答案：}收敛半径 \(R=1\)，收敛区间 \([-1,1]\)，和函数 \[
\boxed{S(x)=x(1+x)\log(1+x)-x^2},\quad-1<x\le1,
\qquad S(-1)=-1.
\] 在 \(|x|<1\) 内，利用 \(1/[n(n-1)]=1/(n-1)-1/n\)， \[
\begin{aligned}
\frac{S(x)}x
&=\sum_{n=2}^\infty\frac{(-1)^nx^n}{n-1}
-\sum_{n=2}^\infty\frac{(-1)^nx^n}{n}\\
&=x\log(1+x)+\log(1+x)-x.
\end{aligned}
\] 在 \(x=0\) 上公式也成立。系数的根值极限给出
\(R=1\)；端点绝对值级数均由 \(\sum_{n\ge2}1/[n(n-1)]\)
控制，故两个端点绝对收敛。由一致收敛延续和函数，\(S(1)=2\log2-1\)，\(S(-1)=-1\)，其中用到
\(u\log u\to0\)。

\subsubsection{4.
高斯余弦积分}\label{ux9ad8ux65afux4f59ux5f26ux79efux5206}

\textbf{答案：}\(\displaystyle\frac{\sqrt\pi}{2}e^{-\alpha^2/4}\)，\(\alpha\in\mathbb R\)。

令 \(J(\alpha)=\int_0^\infty e^{-x^2}\cos(\alpha x)\,dx\)。因为
\(xe^{-x^2}\) 可积，允许在积分号下求导： \[
J'(\alpha)=-\int_0^\infty xe^{-x^2}\sin(\alpha x)\,dx.
\] 利用 \(d(e^{-x^2})=-2xe^{-x^2}dx\) 分部积分，两端边界项为零，得到 \[
\int_0^\infty xe^{-x^2}\sin(\alpha x)\,dx=\frac\alpha2J(\alpha).
\] 故 \(J'=-\alpha J/2\)，结合题给 \(J(0)=\sqrt\pi/2\)，解得上述公式。

\subsection{三、证明题}\label{ux4e09ux8bc1ux660eux9898}

\subsubsection{\texorpdfstring{1. 用 \(\varepsilon\)---\(\delta\)
语言证明有理函数极限}{1. 用 \textbackslash varepsilon---\textbackslash delta 语言证明有理函数极限}}\label{ux7528-varepsilondelta-ux8bedux8a00ux8bc1ux660eux6709ux7406ux51fdux6570ux6781ux9650}

在 \(x\ne1,-1/3\) 时， \[
\frac{x^2-x}{6x^2-4x-2}=\frac{x}{2(3x+1)},\qquad
\left|\frac{x}{2(3x+1)}-\frac18\right|
=\frac{|x-1|}{8|3x+1|}.
\] 给定 \(\varepsilon>0\)，取 \(\delta=\min(1/2,20\varepsilon)\)。当
\(0<|x-1|<\delta\) 时，\(x>1/2\)，故 \(|3x+1|>5/2\)，于是 \[
\left|\frac{x^2-x}{6x^2-4x-2}-\frac18\right|
<\frac{|x-1|}{20}<\varepsilon.
\] 所以所求极限为 \(1/8\)。

\subsubsection{2.
积分不等式的原题条件有误}\label{ux79efux5206ux4e0dux7b49ux5f0fux7684ux539fux9898ux6761ux4ef6ux6709ux8bef}

现有题面给出 \(f(0)=0\)、\(f\) 二阶连续可导，以及 \[
M=\max_{x\in[0,1]}f''(x),\qquad
\left|\int_{-1}^1f(x)\,dx\right|\le\frac M3.
\] \textbf{按这些原始条件，命题不成立。}取 \(f(x)=-x^2\)，则
\(M=-2\)，左侧为 \(2/3\)，右侧为 \(-2/3\)，立即矛盾。

只在 \([0,1]\) 的定义上增加绝对值仍不够。取 \[
f(x)=\begin{cases}x^3,&x<0,\\0,&x\ge0,\end{cases}
\] 则 \(f\in C^2[-1,1]\)，\(f(0)=0\)，\(\max_{[0,1]}|f''|=0\)，但
\(|\int_{-1}^1f|=1/4>0\)。

\textbf{助手建议的正确条件版本：}将 \(M\) 定义为
\(\max_{x\in[-1,1]}|f''(x)|\)。这同时修正了控制区间和绝对值，未将修正条件冒充原文。此时由
Taylor 公式 \[
f(x)=f'(0)x+R(x),\qquad |R(x)|\le\frac M2x^2.
\] 奇函数项 \(f'(0)x\) 的对称区间积分为零，于是 \[
\left|\int_{-1}^1f(x)\,dx\right|
\le\int_{-1}^1|R(x)|\,dx
\le\frac M2\int_{-1}^1x^2\,dx=\frac M3.
\] 取 \(f(x)=cx^2\) 可达等号，因此该正确版本的常数 \(1/3\) 最优。

\subsubsection{3.
验证弦振动方程}\label{ux9a8cux8bc1ux5f26ux632fux52a8ux65b9ux7a0b}

作经典二阶导数验证需要
\(a\ne0\)、\(f\in C^2\)、\(F\in C^1\)；原题没有列出光滑性，应将其视为公式所需的适用条件。设
\(p=x-at,q=x+at\)，由链式法则和变上限积分求导法则， \[
\begin{aligned}
u_x&=\frac12[f'(p)+f'(q)]+\frac1{2a}[F(q)-F(p)],\\
u_{xx}&=\frac12[f''(p)+f''(q)]+\frac1{2a}[F'(q)-F'(p)],\\
u_t&=\frac a2[-f'(p)+f'(q)]+\frac12[F(q)+F(p)],\\
u_{tt}&=\frac{a^2}2[f''(p)+f''(q)]+\frac a2[F'(q)-F'(p)].
\end{aligned}
\] 逐项比较即得 \(u_{tt}=a^2u_{xx}\)。公式还给出
\(u(x,0)=f(x)\)、\(u_t(x,0)=F(x)\)。

\subsubsection{4.
三角函数和的一致收敛}\label{ux4e09ux89d2ux51fdux6570ux548cux7684ux4e00ux81f4ux6536ux655b}

将题中的有限和理解为函数列 \[
S_n(x)=\frac1n\sum_{k=1}^n\cos\left(x+\frac kn\right).
\] 其极限为 \[
S(x)=\int_0^1\cos(x+t)\,dt=\sin(x+1)-\sin x.
\] 利用 \(|\cos u-\cos v|\le|u-v|\)，对任意实数 \(x\)， \[
\begin{aligned}
|S_n(x)-S(x)|
&\le\sum_{k=1}^n\int_{(k-1)/n}^{k/n}
|\cos(x+k/n)-\cos(x+t)|\,dt\\
&\le\sum_{k=1}^n\int_{(k-1)/n}^{k/n}(k/n-t)\,dt
=\frac1{2n}\to0.
\end{aligned}
\] 上界与 \(x\) 无关，因而在 \([a,b]\)
上一致收敛；实际上在整个实数轴上也一致收敛。

\subsubsection{5.
一致连续函数的一致极限}\label{ux4e00ux81f4ux8fdeux7eedux51fdux6570ux7684ux4e00ux81f4ux6781ux9650}

给定 \(\varepsilon>0\)。由一致收敛，存在 \(N\) 使所有 \(x\in I\) 均有
\(|f_N(x)-f(x)|<\varepsilon/3\)。由固定函数 \(f_N\) 的一致连续性，存在
\(\delta>0\)，当 \(x,y\in I\) 且 \(|x-y|<\delta\)
时，\(|f_N(x)-f_N(y)|<\varepsilon/3\)。于是 \[
|f(x)-f(y)|\le|f(x)-f_N(x)|+|f_N(x)-f_N(y)|+|f_N(y)-f(y)|<\varepsilon.
\] 因此 \(f\) 在 \(I\) 上一致连续。不同 \(f_n\) 的一致连续模数无需统一。

\subsection{四、讨论题}\label{ux56dbux8ba8ux8bbaux9898}

\subsubsection{1.
椭球切平面与坐标轴围成的四面体}\label{ux692dux7403ux5207ux5e73ux9762ux4e0eux5750ux6807ux8f74ux56f4ux6210ux7684ux56dbux9762ux4f53}

\textbf{按题面所问 \(V_{\max}\)，体积没有有限上界。}通常假设半轴参数
\(\alpha,\beta,\gamma>0\)。在椭球上取切点
\((x_0,y_0,z_0)\)，三个坐标均非零，切平面为 \[
\frac{x_0x}{\alpha^2}+\frac{y_0y}{\beta^2}+\frac{z_0z}{\gamma^2}=1.
\] 其三个截距是 \(\alpha^2/x_0,\beta^2/y_0,\gamma^2/z_0\)，所以 \[
V=\frac{\alpha^2\beta^2\gamma^2}{6|x_0y_0z_0|}.
\] 在第一卦限内取 \[
x_0=\alpha\varepsilon,\quad
y_0=\beta\sqrt{\frac{1-\varepsilon^2}{2}},\quad
z_0=\gamma\sqrt{\frac{1-\varepsilon^2}{2}},\qquad0<\varepsilon<1.
\] 这些点均在椭球上，所围四面体均有限，但 \[
V=\frac{\alpha\beta\gamma}{3\varepsilon(1-\varepsilon^2)}\to+\infty
\quad(\varepsilon\downarrow0).
\] 因此即使限制切点与截距均在第一卦限，最大值仍不存在。

\textbf{若原题实际要问最小值：}令
\(u=|x_0|/\alpha,v=|y_0|/\beta,w=|z_0|/\gamma\)，则
\(u^2+v^2+w^2=1\)。由均值不等式，\(uvw\le1/(3\sqrt3)\)，故 \[
\boxed{V_{\min}=\frac{\sqrt3}{2}\alpha\beta\gamma},
\] 等号在 \(u=v=w=1/\sqrt3\) 时成立。这是另一个正确问题的答案，不是原文
\(V_{\max}\) 的答案。

\subsubsection{2.
有理数与无理数分段函数}\label{ux6709ux7406ux6570ux4e0eux65e0ux7406ux6570ux5206ux6bb5ux51fdux6570}

\[
f(x)=\begin{cases}x(1-x),&x\in\mathbb Q,\\x(1+x),&x\notin\mathbb Q.\end{cases}
\] \textbf{结论：}只在 \(x=0\) 连续，只在 \(x=0\) 可微，且 \(f'(0)=1\)。

有理数、无理数均稠密。在点 \(a\)
连续必须且只须两条多项式分支在该点一致，即 \(a-a^2=a+a^2\)，解得
\(a=0\)。在零点有
\(|f(x)|\le|x|+x^2\to0=f(0)\)，所以确实连续。非零点不连续，因而不可微。

在零点，差商为 \[
\frac{f(h)-f(0)}h=\begin{cases}1-h,&h\in\mathbb Q\setminus\{0\},\\1+h,&h\notin\mathbb Q,\end{cases}
\] 两种情况下均趋于 \(1\)，故 \(f'(0)=1\)。

\subsection{核验说明}\label{ux6838ux9a8cux8bf4ux660e}

题面已交叉读取远程OCR与本地逐页OCR，其中一6的被积式由本地OCR补全。使用
SymPy
对积分、球面积分的径向与角向因子、二2的交叉偏导和势函数、二3的级数系数以及三2反例作符号检查；检查脚本为
\texttt{.agent/hdu-ocr/check\_math2024\_2025.py}。一7题干仍缺失；二2曲线条件、三2的
\(M\) 定义、四1的最大值问题均保留明确说明。
\end{document}
